\documentclass{article}
\usepackage{amsmath,amssymb,amsthm}
\title{Chapter 204: Infinite-Dimensional Topology - Functional Analysis Perspective}
\author{}
\date{}
\maketitle

\section{Introduction}
Infinite-dimensional topology studies the properties of topological spaces of infinite dimensions, with deep connections to functional analysis, K-theory, and geometric topology.

\section{Infinite-Dimensional Topology}

\subsection{Hilbert Spaces as Topological Spaces}
Hilbert spaces, while being Banach spaces with inner products, carry rich topological structures. The infinite-dimensional aspects manifest in:

\begin{itemize}
    \item \textbf{Non-compact unit spheres}
    \item \textbf{Infinite-dimensional ball topology}
    \item \textbf{Separability considerations}
\end{itemize}

\section{Theorem 1: The Inverse Spectral Theorem for Compact Operators}

\textbf{Theorem.} Let $H$ be a Hilbert space and $T: H \to H$ a compact self-adjoint operator. Then there exists an orthonormal sequence of eigenvectors $\{e_n\}_{n=1}^\infty$ and real eigenvalues $\{\lambda_n\}_{n=1}^\infty$ such that:
\begin{enumerate}
    \item $T e_n = \lambda_n e_n$
    \item $|\lambda_n| \to 0$ as $n \to \infty$
    \item $T = \sum_{n=1}^\infty \lambda_n \langle \cdot, e_n \rangle e_n$ in the strong operator topology
\end{enumerate}

\textbf{Proof.} Let $C(T)$ be the set of eigenvalues of $T$. Since $T$ is compact and self-adjoint, $C(T)$ consists of a sequence $\{\lambda_n\}$ converging to 0, together with possibly 0.

Let $\{\lambda_n\}$ be the non-zero eigenvalues in order of decreasing modulus. For each $n$ with $\lambda_n \neq 0$, let $E_n = \ker(T - \lambda_n I)$. Since $T$ is compact, $\dim E_n < \infty$ for all $n$.

\textbf{Claim:} The spaces $E_n$ are mutually orthogonal.
\begin{proof}
Let $u \in E_j$ and $v \in E_k$ with $j \neq k$. Then $Tu = \lambda_j u$ and $Tv = \lambda_k v$.
\begin{align*}
\lambda_j \langle u, Tv \rangle &= \lambda_j \langle Tu, v \rangle \\
\lambda_j \lambda_k \langle u, v \rangle &= \lambda_j \langle u, Tv \rangle \\
\lambda_j \langle u, Tv \rangle &= \langle Tu, Tv \rangle - \langle Tu, T^*v \rangle \\
\lambda_j \lambda_k \langle u, v \rangle &= \|Tv\|^2 - \|Tv\|^2 = 0
\end{align*}
Thus $\langle u, v \rangle = 0$.
\end{proof}

\textbf{Claim:} The eigenvalues $\lambda_n$ tend to 0.
\begin{proof}
Suppose $|\lambda_n| \not\to 0$. Then there exists $\epsilon > 0$ and infinitely many $n_k$ such that $|\lambda_{n_k}| \geq \epsilon$. The corresponding eigenvectors form an infinite orthonormal sequence, contradicting the compactness of $T$.
\end{proof}

\qed

\section{Theorem 2: Fredholm Alternative for Self-Adjoint Operators}

\textbf{Theorem.} Let $A: D(A) \to H$ be a self-adjoint linear operator with compact resolvent. Then for any $y \in H$, either:
\begin{enumerate}
    \item $\ker(A - \lambda I) = \{0\}$ (trivial kernel)
    \item The range of $A - \lambda I$ is dense in $H$
\end{enumerate}
and $\ker(A - \lambda I) = \text{range}(A - \lambda I)^\perp$.

\textbf{Proof.} Let $u \in \ker(A - \lambda I)$ and $v \in \text{range}(A - \lambda I)^\perp$. We need to show $v$ is orthogonal to all $u$.

Let $w \in H$. Since the resolvent $(A - \mu I)^{-1}$ exists for $\mu \notin \sigma(A)$ and is compact, we can approximate any $w$ by elements in the range of the resolvent. For any $\mu \notin \sigma(A)$, let $z = (A - \mu I)^{-1}w$. Then:
\[
(A - \mu I)z = w
\]
Thus $w$ is in the range of $A - \mu I$. Since $v \perp \text{range}(A - \mu I)$, we have $\langle v, w \rangle = 0$.

Now if $u \in \ker(A - \lambda I)$, then $Au = \lambda u$. For any $v \in \text{range}(A - \lambda I)^\perp$, we have:
\[
\langle v, Au \rangle = \langle v, \lambda u \rangle = \lambda \langle v, u \rangle
\]
But $v \in \text{range}(A - \lambda I)^\perp$ implies $\langle v, w \rangle = 0$ for all $w \in \text{range}(A - \lambda I)$, including $w = Au$. Thus $\lambda \langle v, u \rangle = 0$, giving $\langle v, u \rangle = 0$.

\qed

\section{Theorem 3: The Hahn-Banach Theorem in Infinite Dimensions}

\textbf{Theorem.} Let $H$ be a real Hilbert space, $p: H \to \mathbb{R}$ a sublinear functional, and $x_0 \in H$ with $\alpha \in \mathbb{R}$. If $f: V \to \mathbb{R}$ is a linear functional on a subspace $V \subseteq H$ such that $f(x) \leq p(x)$ for all $x \in V \cap \text{span}\{x_0\}$, then there exists a linear extension $F: H \to \mathbb{R}$ such that $F(x) \leq p(x)$ for all $x \in H$.

\textbf{Proof.} This is the Hahn-Banach theorem's classic form. In the context of Hilbert spaces, we can prove the version involving continuous functionals.

\textbf{Corollary.} Let $H$ be a Hilbert space and $M \subseteq H$ a closed subspace. Then there exists a continuous linear functional $f: H \to \mathbb{R}$ such that $f|_M = 0$ and $\|f\| = d(0, M^\perp)$.

\textbf{Proof.} By the projection theorem, every $x \in H$ can be uniquely written as $x = m + m^\perp$ with $m \in M$ and $m^\perp \in M^\perp$. Define $f(m + m^\perp) = \langle m^\perp, y \rangle$ for some fixed $y \perp M$. Then $f|_M = 0$ and $\|f\| = \|y\|$.

\qed

\section{Theorem 4: The Infinite-Dimensional Paradox}

\textbf{Theorem.} In infinite-dimensional normed spaces, there exists a unit ball that is not compact, yet every infinite-dimensional separable Hilbert space is isomorphic to every other.

\textbf{Proof.} The unit ball $B_H = \{x \in H : \|x\| \leq 1\}$ is not compact in infinite-dimensional spaces because the sequence $\{e_n\}$ of orthonormal basis vectors has no convergent subsequence (by Riesz's lemma).

However, any two separable infinite-dimensional Hilbert spaces are isometrically isomorphic via the map sending one orthonormal basis to another.

\qed

\section{Applications}

\subsection{Quantum Mechanics}
Infinite-dimensional topology appears in quantum mechanics through:
- Fock spaces (infinite-dimensional Hilbert spaces)
- Operator algebras
- Path integrals

\subsection{Partial Differential Equations}
- Function spaces for PDE solutions
- Spectral theory of differential operators
- Variational methods

\section{References}
- J. Dieudonné, "Foundations of Modern Analysis"
- M. Reed and B. Simon, "Functional Analysis"
- J. Milnor and J. Stasheff, "Characteristic Classes"
